%
% QNFO STANDARD PUBLICATION TEMPLATE
% Based on the official Springer Nature LaTeX Template (sn-jnl.cls, v3.1, December 2024)
% Downloaded from: https://www.springernature.com/gp/authors/campaigns/latex-author-support
% This is the CURRENT, OFFICIAL Springer Nature template -- NOT the legacy svjour3.cls
% (svjour3/svjour.cls, distributed via CTAN's "springer" package, was retired by
% Springer Nature in favor of the unified sn-jnl.cls across all SN journals in ~2019-2021.
% Foundations of Physics accepts the current SN template; do not chase svjour3 for new work.)
%
% BUILD REQUIREMENTS (verified working 2026-07-25 with TeX Live 2025 on Windows):
%   1. Place sn-jnl.cls in the SAME directory as this .tex file.
%   2. Copy the correct .bst file (e.g. sn-mathphys-num.bst) from bst/ into the
%      SAME directory as this .tex file -- bibtex does NOT search subdirectories
%      by default and will fail with "I couldn't open style file ..." otherwise.
%      (This is a genuine kaizen finding from the 2026-07-25 template validation
%      session -- see research/SKILL.md Anti-Patterns table.)
%   3. Build sequence: pdflatex -> bibtex -> pdflatex -> pdflatex (4 passes total,
%      standard LaTeX+BibTeX convergence pattern).
%
\documentclass[pdflatex,sn-mathphys-num]{sn-jnl}% Math and Physical Sciences Numbered Reference Style
%% Alternative class options for other QNFO paper types:
%%   sn-mathphys-ay   -- Author-Year citation style instead of numbered
%%   sn-basic         -- Generic Springer Nature style (non-physics)
%% See templates/springer-nature-latex/user-manual.pdf for the full option list.

\usepackage{graphicx}%
\usepackage{amsmath,amssymb,amsfonts}%
\usepackage{amsthm}%
\usepackage{mathrsfs}%
\usepackage[title]{appendix}%
\usepackage{xcolor}%
\usepackage{textcomp}%
\usepackage{booktabs}%

% --- QNFO convention: bra-ket / physics notation support ---
% Uncomment if the paper uses \ket{}, \bra{}, \braket{} directly (rather than
% relying on unicode-latex-preprocess.py's \langle/\rangle fallback):
%\usepackage{braket}

\theoremstyle{thmstyleone}%
\newtheorem{theorem}{Theorem}%
\newtheorem{proposition}[theorem]{Proposition}%
\newtheorem{lemma}[theorem]{Lemma}%
\newtheorem{corollary}[theorem]{Corollary}%

\theoremstyle{thmstyletwo}%
\newtheorem{example}{Example}%
\newtheorem{remark}{Remark}%

\theoremstyle{thmstylethree}%
\newtheorem{definition}{Definition}%

\raggedbottom

\begin{document}

\title[Short Running Title]{Full Paper Title Goes Here}

% --- QNFO convention: single-author attribution ---
\author*[1]{\fnm{Rowan Brad} \sur{Quni-Gudzinas}}\email{author@qnfo.org}

\affil*[1]{\orgname{QNFO}, \orgaddress{\country{}}}

% --- QNFO convention: DOI + license disclosed at top, not just in Declarations ---
% Add DOI once Zenodo deposit is minted (Phase 5 of research/SKILL.md).
% Placeholder during drafting: DOI: 10.5281/zenodo.PENDING

\abstract{Replace with a 150--250 word abstract per Foundations of Physics
(Springer Nature) submission guidelines. The abstract should not contain
undefined abbreviations or unspecified references. No subheadings, equations,
or citations unless the target journal explicitly permits them.}

\keywords{keyword1, keyword2, keyword3, keyword4}

\maketitle

% ======================================================================
% 1. INTRODUCTION
% ======================================================================
\section{Introduction}\label{sec1}

Replace with introduction content. Cite via \verb+\cite{key}+ (numbered
style under sn-mathphys-num).

% ======================================================================
% 2. [SECTION TITLE]
% ======================================================================
\section{Section Title}\label{sec2}

Body content.

% ======================================================================
% N. CONCLUSION
% ======================================================================
\section{Conclusion}\label{secN}

Conclusion content.

% ======================================================================
% DECLARATIONS (QNFO MANDATORY BLOCK -- do not omit any subsection)
% ======================================================================
\section*{Declarations}

\subsection*{Funding}
This research received no external funding and was conducted as part of the QNFO foundational physics program.

\subsection*{Conflicts of Interest}
The author declares no conflicts of interest.

\subsection*{Ethics Approval and Consent to Participate}
Not applicable.

\subsection*{Consent for Publication}
Not applicable.

\subsection*{Author Contributions}
R.B.Q.-G. conceived the argument, verified all claims against primary sources, and takes full responsibility for the content of this manuscript.

\subsection*{Data Availability}
No experimental data were generated or analyzed in this study. All cited literature is publicly available via their respective publishers or preprint servers.

\subsection*{Materials Availability}
Not applicable.

\subsection*{Code Availability}
Not applicable, or: all analysis code is publicly available at \url{https://github.com/rwnq8/REPO-NAME}.

\subsection*{Use of Artificial Intelligence}
This manuscript was drafted with the assistance of the DeepChat large language
model, which conducted literature search support, composed initial text
drafts, and formatted references. Per the Springer Nature policy on
AI-assisted authoring \cite{SpringerAI2024}, this use is disclosed here in the
Declarations section. All substantive claims, arguments, and citations were
verified by the author against primary sources. The author takes full
accountability for the final version of the text and attests that all
citations reference genuine, published works. No large language model is
listed as an author.

% ======================================================================
% BIBLIOGRAPHY
% ======================================================================
\bibliography{refs}

\end{document}
